\hypertarget{appendix-b0732b22e7d6}{}
\section*{Additive energy and approximate structure}
\section{Additive energy and approximate structure}\label{additive-energy-and-approximate-structure}

\textbf{Curriculum depth:} Developed · \textbf{First seen:} University courses · \textbf{Latest development:} University courses

\subsection{Definition and mechanism}\label{definition-and-mechanism}

Additive energy measures how often two pairs from a set have the same sum: \(E(A)=|\{(a,b,c,d)\in A^4:a+b=c+d\}|\). A random-like set has relatively few repeated sums; a set with many additive coincidences often contains a structured subset. The Balog--Szemerédi--Gowers principle turns high energy into a large subset with small doubling, while Ruzsa-style inequalities and Freiman-type theorems describe how small doubling constrains approximate additive structure. In finite fields, sum-product results show that a set cannot remain simultaneously additively and multiplicatively concentrated unless it has special algebraic structure, such as subfield behavior. These implications are quantitative and depend on group or field hypotheses, subset size, and constants. High energy does not mean the entire set is a subgroup, and a small-doubling conclusion may apply only after discarding many elements. Energy is a counting diagnostic that guides structure theorems, not a complete classification by itself.

\subsection{Prerequisites}\label{prerequisites}

\begin{itemize}
\tightlist
\item
  \href{}{Vector spaces subspaces and span (02-Foundations-II, appendix-173b52c063ae)}
\item
  \href{}{Basis coordinates and dimension (01-Foundations-I, appendix-ce0341f7ea0d)}
\item
  \href{}{Scalars fields and real-number models (02-Foundations-II, appendix-52051ad550fe)}
\item
  \hyperref[appendix-3c73f7862157]{Exceptional projections and dimension}
\end{itemize}

\subsection{Used by}\label{used-by}

\begin{itemize}
\tightlist
\item
  \href{Continuous\%20Notes/University\%20-\%20Worlds,\%20Images\%20and\%20Inference\%20-\%20Graphics,\%20Vision\%20and\%20Photography\#Additive\%20structure\%20explains\%20repeated\%20coincidences}{Additive structure explains repeated coincidences}
\item
  \href{Continuous\%20Notes/University\%20-\%20Worlds,\%20Images\%20and\%20Inference\%20-\%20Graphics,\%20Vision\%20and\%20Photography\#Dynamics\%20connect\%20projection\%20growth\%20to\%20rich\%20tubes}{Dynamics connect projection growth to rich tubes}
\end{itemize}

\subsection{Full teaching explanation}\label{full-teaching-explanation}

\begin{itemize}
\tightlist
\item
  \href{Continuous\%20Notes/University\%20-\%20Worlds,\%20Images\%20and\%20Inference\%20-\%20Graphics,\%20Vision\%20and\%20Photography\#Additive\%20structure\%20explains\%20repeated\%20coincidences}{Additive structure explains repeated coincidences}
\item
  \href{Continuous\%20Notes/University\%20-\%20Worlds,\%20Images\%20and\%20Inference\%20-\%20Graphics,\%20Vision\%20and\%20Photography\#Dynamics\%20connect\%20projection\%20growth\%20to\%20rich\%20tubes}{Dynamics connect projection growth to rich tubes}
\end{itemize}

\subsection{Recall and next depth}\label{recall-and-next-depth}

\begin{itemize}
\tightlist
\item
  Count one solution to a+b=c+d and explain what high energy suggests about a subset.
\item
  Why does a BSG conclusion not say the entire set is a subgroup?
\end{itemize}

\subsection{Source and navigation}\label{source-and-navigation}

Teaching narrative: \href{Continuous\%20Notes/University\%20-\%20Worlds,\%20Images\%20and\%20Inference\%20-\%20Graphics,\%20Vision\%20and\%20Photography\#Additive\%20structure\%20explains\%20repeated\%20coincidences}{open the main explanation}. Browse the \href{}{University graphics and vision index (02-Foundations-II, appendix-a6d556a4ef50 § University graphics vision and photography - independent series)}. Additional chapter connections: \href{Continuous\%20Notes/Computing\%20Foundations\%20-\%20From\%20Bits\%20to\%20C++\%20and\%20Systems\%20APIs\#A\%20linear\%20equation\%20is\%20a\%20measurement,\%20and\%20several\%20equations\%20are\%20geometry}{A linear equation is a measurement, and several equations are geometry}; \href{Continuous\%20Notes/University\%20-\%20Worlds,\%20Images\%20and\%20Inference\%20-\%20Graphics,\%20Vision\%20and\%20Photography\#Sampling\%20continuous\%20signals\%20and\%20image\%20formation}{Sampling continuous signals and image formation}; \href{Continuous\%20Notes/Computing\%20Foundations\%20-\%20From\%20Bits\%20to\%20C++\%20and\%20Systems\%20APIs\#The\%20reachable\%20movements\%20form\%20a\%20vector\%20space}{The reachable movements form a vector space}.

Caption/source records: \href{Continuous\%20Notes/University\%20-\%20Worlds,\%20Images\%20and\%20Inference\%20-\%20Graphics,\%20Vision\%20and\%20Photography\#P2-09\%20-\%20Lecture\%2009:\%20Reflections\%20on\%20the\%20Szemeredi-Trotter\%20Theorem}{P2-09}, \href{Continuous\%20Notes/University\%20-\%20Worlds,\%20Images\%20and\%20Inference\%20-\%20Graphics,\%20Vision\%20and\%20Photography\#P2-10\%20-\%20Lecture\%2010:\%20Sum-Product\%20Theory}{P2-10}, \href{Continuous\%20Notes/University\%20-\%20Worlds,\%20Images\%20and\%20Inference\%20-\%20Graphics,\%20Vision\%20and\%20Photography\#P2-11\%20-\%20Lecture\%2011:\%20Contagious\%20Structure\%20in\%20Projection\%20Theory}{P2-11}, \href{Continuous\%20Notes/University\%20-\%20Worlds,\%20Images\%20and\%20Inference\%20-\%20Graphics,\%20Vision\%20and\%20Photography\#P2-16\%20-\%20Lecture\%2016:\%20The\%20Bourgain\%20Projection\%20Theorem,\%20Part\%203}{P2-16}.
